% !TEX root = ../main.tex
\section{Archived Extended Baselines}
\label{ch:appendix-archive}

During exploratory analysis, additional PageRank variants were scored on the same matched cohort and proxy families. They are archived here for completeness and future work. They are \emph{not} part of the primary EndorseRank--AWP research questions and should not be read as peer claims in the main narrative.

\dissertationSubheading[sec:seven-method-archive]{Seven-method alignment matrix}

Exploratory seven-method matrices remain in the evaluation archive (LP-PR, CW-AWP, LF-PR, and RiskProp-style variants). They are not drawn here. EndorseRank remains the allowance-family leader among social methods; AWP (and closely related transfer variants) lead on transfer.

Interpretation guidelines:

\begin{itemize}
\item Methods built on transfer-like edges cluster near AWP on transfer and Sybil families.
\item Methods with very small edge counts can produce unstable or degenerate rankings; their runtimes may be low simply because $m$ is tiny.
\item RiskProp-style rows that mirror AWP numerically indicate shared transfer structure, not independent discovery.
\end{itemize}

\dissertationSubheading[sec:six-aave-archive]{Six-method Aave-oriented matrix}

An Aave-oriented preset remains in the evaluation archive. Several lending-event graphs are sparse on Arbitrum within the study window; delegation-based rows may be degenerate when \texttt{BorrowAllowanceDelegated} events were not extracted (Appendix~\ref{ch:appendix-eval}).

The Aave-oriented preset remains in the evaluation archive and is not drawn here.

These matrices motivate future work on credit-delegation edges and cross-protocol lending labels, but they do not alter Claims C1--C4.

\dissertationSubheading[sec:why-archived]{Why the main narrative uses two methods}

Focusing this research on EndorseRank and AWP serves two goals:

\begin{enumerate}
\item Clarity of contribution: The novel construct is allowance-based endorsement. Comparing it to the established transfer baseline\cend{3} isolates the research question.
\item Avoiding baseline sprawl: Additional variants that share transfer edges with AWP add rows without adding independent constructs.
\item Avoiding circular references: The outcome-native GF-PR reference used in the research proposal is built from the same GMX events as the trading-success proxies, so any alignment between them is construct overlap, not validation (Section~\ref{sub:gf-pr}); its three-method matrix is archived alongside the seven-method and six-Aave presets.
\end{enumerate}

Researchers who need broader baselines can regenerate the archive tables from the open evaluation pipeline using the seven-method and six-Aave presets.

\dissertationSubheading[sec:archive-repro]{Reproducing archive tables}

Archive tables are regenerated by the evaluation pipeline's table export with seven-method and six-Aave presets. CSV copies are recorded alongside the LaTeX fragments for spreadsheet inspection.

\dissertationSubheading[sec:archive-method-notes]{Notes on archived method identifiers}

The following identifiers appear in archive tables and CSV exports. They document exploratory implementations in the evaluation codebase.

\begin{description}
\item[LP-PR] A transfer-graph variant with alternative loss-penalized or liquidity-oriented weighting explored during development. Empirically it tracks AWP closely on transfer families, indicating shared payment-flow structure.
\item[CW-AWP] A counterparty-weighted AWP variant that down-weights repeated transfers between the same pair. It slightly reduces transfer $\tau$ relative to AWP while remaining in the same construct family.
\item[LF-PR] A low-frequency or lightly filtered graph with very small edge counts and near-zero runtimes; rankings are not comparable to full AWP and are retained only to document a failed sparse baseline.
\item[RiskProp] A risk-propagation style ranking on transfer-like edges inspired by Lin et al.\cend{12}. On this cohort it numerically mirrors AWP on several families, suggesting that the implemented risk signal did not diverge strongly from transfer centrality under the study's parameterization.
\item[LiqCall-PR / Borrow-PR / Repay-PR] Aave-oriented wallet-to-wallet graphs built from liquidation calls, borrows, and repays when those events are present. On Arbitrum in the study window, edge counts are small (tens to hundreds), so alignment is weak and runtimes are negligible.
\item[Delegation-PR] Intended to use credit-delegation allowances; with zero delegation edges extracted, scores are degenerate and $\tau$ is not interpretable.
\end{description}

These notes exist so that future researchers can decide which variants deserve full re-implementation. Archive rows are not peer-reviewed baselines.

\dissertationSubheading[sec:archive-decision-log]{Decision log: what stayed in the main text}

The decision to keep only EndorseRank and AWP in Chapters~\ref{ch:results}--\ref{ch:discussion} was made after inspecting the seven-method and six-Aave matrices. Criteria were:

\begin{enumerate}
\item Construct independence: A baseline must implement a distinct edge semantics, not a minor reweighting of transfers.
\item Data adequacy: Edge counts must be large enough for stable PageRank on the matched cohort.
\item Narrative focus: This research's contribution is allowance-based endorsement, not a survey of all possible PageRank variants.
\end{enumerate}

Under these criteria, EndorseRank (allowance) and AWP (transfer) form a minimal sufficient set. Archive tables remain available for readers who want the seven-method and six-Aave matrices without mixing them into the primary claims.

